\documentclass{beamer}
\usetheme{Madrid}
\usecolortheme{whale}

\usepackage[utf8]{inputenc}
\usepackage[spanish]{babel}
\usepackage{amsmath}
\usepackage{tikz}
\usetikzlibrary{arrows, shapes.gates.logic.US, calc}

% Configuración de elementos para diagramas de Op-Amps
\tikzset{
    opamp/.style={
        triangle,                 
        draw,                     
        minimum size=2cm,         
        regular polygon sides=3,  
        shape border rotate=-90,  
        thick
    }
}

\title[Amplificadores Operacionales]{Configuraciones Básicas de Amplificadores Operacionales}
\author{Electrónica Analógica}
\date{\today}

\begin{document}

% Diapositiva 1: Portada
\begin{frame}
    \titlepage
\end{frame}

% Diapositiva 2: Introducción
\begin{frame}{Introducción al Amplificador Operacional}
    \begin{block}{¿Qué es un Op-Amp?}
        Es un dispositivo amplificador de alta ganancia con entrada diferencial y, por lo general, una salida única. Es el bloque constructivo fundamental de la electrónica analógica.
    \end{block}
    \vspace{0.3cm}
    \textbf{Características del Op-Amp Ideal:}
    \begin{itemize}
        \item Ganancia de lazo abierto infinita ($A_{OL} = \infty$)
        \item Impedancia de entrada infinita ($Z_{in} = \infty$)
        \item Impedancia de salida cero ($Z_{out} = 0$)
        \item Ancho de banda infinito
    \end{itemize}
\end{frame}

% Diapositiva 3: Seguidor de Voltaje
\begin{frame}{1. Seguidor de Voltaje (Buffer)}
    \begin{columns}
        \begin{column}{0.5\textwidth}
            \textbf{Ecuación:}
            \[ V_{out} = V_{in} \]
            \vspace{0.2cm}
            \textbf{Características:}
            \begin{itemize}
                \item Ganancia de voltaje unitaria ($A_v = 1$).
                \item Alta impedancia de entrada y baja impedancia de salida.
                \item Se utiliza para acoplamiento de impedancias (aislar etapas).
            \end{itemize}
        \end{column}
        \begin{column}{0.5\textwidth}
            \begin{centering}
            \begin{tikzpicture}[scale=0.7, transform shape]
                \draw (0,0) node[opamp] (opamp) {};
                \draw (opamp.corner 3) node[left, yshift=0.4cm] {$+$} node[left, yshift=-0.4cm] {$-$};
                \draw (opamp.corner 3) ++(-1.5,0.4) node[left] {$V_{in}$} -- ++(1,0) -- +(0.5,0);
                \draw (opamp.corner 1) -- ++(1,0) node[right] {$V_{out}$} coordinate (out);
                \draw (opamp.corner 3) ++(0,-0.4) -- ++(0,-0.8) coordinate (fb) -- ($(out)+(0,-1.2)$) -- (out);
            \end{tikzpicture}
            \end{centering}
        \end{column}
    \end{columns}
\end{frame}

% Diapositiva 4: Amplificador Inversor
\begin{frame}{2. Amplificador Inversor}
    \begin{columns}
        \begin{column}{0.5\textwidth}
            \textbf{Ecuación:}
            \[ V_{out} = -\left(\frac{R_f}{R_{in}}\right) V_{in} \]
            \vspace{0.2cm}
            \textbf{Características:}
            \begin{itemize}
                \item La salida está desfasada $180^\circ$ respecto a la entrada.
                \item El pin no inversor ($+$) está conectado a tierra (Tierra Virtual en pin $-$).
                \item La ganancia puede ser menor, igual o mayor a 1.
            \end{itemize}
        \end{column}
        \begin{column}{0.5\textwidth}
            \begin{centering}
            \begin{tikzpicture}[scale=0.7, transform shape]
                \draw (0,0) node[opamp] (opamp) {};
                \draw (opamp.corner 3) node[left, yshift=0.4cm] {$+$} node[left, yshift=-0.4cm] {$-$};
                % Entrada inversora
                \draw (opamp.corner 3) ++(0,-0.4) -- ++(-0.5,0) to[R, l=$R_{in}$] ++(-1.5,0) node[left] {$V_{in}$};
                % Lazo de realimentacion
                \draw (opamp.corner 3) ++(-0.2,-0.4) coordinate (node1) -- ++(0,1.2) to[R, l=$R_f$] ++(1.8,0) -- ++(0,-1.2) coordinate (node2);
                \draw (opamp.corner 1) -- (node2) -- ++(0.5,0) node[right] {$V_{out}$};
                % Entrada no inversora a tierra
                \draw (opamp.corner 3) ++(0,0.4) -- ++(-0.5,0) -- ++(0,-0.3) node[ground] {};
                \draw (-0.5,0.1) -- ++(0,-0.2) -- ++(-0.2,0) ++(0.4,0) -- ++(-0.4,0); % Simbolo tierra simple
            \end{tikzpicture}
            \end{centering}
        \end{column}
    \end{columns}
\end{frame}

% Diapositiva 5: Amplificador No Inversor
\begin{frame}{3. Amplificador No Inversor}
    \begin{columns}
        \begin{column}{0.5\textwidth}
            \textbf{Ecuación:}
            \[ V_{out} = \left(1 + \frac{R_f}{R_1}\right) V_{in} \]
            \vspace{0.2cm}
            \textbf{Características:}
            \begin{itemize}
                \item La salida está en fase con la señal de entrada.
                \item La ganancia de voltaje siempre es mayor o igual a 1.
                \item Ofrece una impedancia de entrada extremadamente alta.
            \end{itemize}
        \end{column}
        \begin{column}{0.5\textwidth}
            \begin{centering}
            \begin{tikzpicture}[scale=0.7, transform shape]
                \draw (0,0) node[opamp] (opamp) {};
                \draw (opamp.corner 3) node[left, yshift=0.4cm] {$+$} node[left, yshift=-0.4cm] {$-$};
                % Entrada no inversora
                \draw (opamp.corner 3) ++(-0.5,0.4) node[left] {$V_{in}$} -- ++(0.5,0);
                % Realimentacion
                \draw (opamp.corner 3) ++(-0.2,-0.4) coordinate (n1) -- ++(0,-1.2) to[R, l=$R_f$] ++(1.8,0) -- ++(0,1.2) coordinate (n2);
                \draw (opamp.corner 1) -- ++(1,0) node[right] {$V_{out}$};
                % R1 a tierra
                \draw (opamp.corner 3) ++(-0.2,-0.4) -- ++(-1,0) to[R, l=$R_1$] ++(-1,0) -- ++(0,-0.3);
                \draw (-2.4,-0.7) -- ++(0,-0.2) -- ++(-0.2,0) ++(0.4,0); % Tierra
            \end{tikzpicture}
            \end{centering}
        \end{column}
    \end{columns}
\end{frame}

% Diapositiva 6: Amplificador Sumador
\begin{frame}{4. Amplificador Sumador Inversor}
    \begin{columns}
        \begin{column}{0.5\textwidth}
            \textbf{Ecuación:}
            \[ V_{out} = -R_f \left( \frac{V_1}{R_1} + \frac{V_2}{R_2} \right) \]
            \vspace{0.2cm}
            \textbf{Si $R_1 = R_2 = R_f$:}
            \[ V_{out} = -(V_1 + V_2) \]
            \textbf{Características:}
            \begin{itemize}
                \item Suma algebraicamente múltiples señales de entrada.
                \item Cada entrada puede escalarse de forma independiente variando sus resistencias.
            \end{itemize}
        \end{column}
        \begin{column}{0.5\textwidth}
            \begin{centering}
            \begin{tikzpicture}[scale=0.6, transform shape]
                \draw (0,0) node[opamp] (opamp) {};
                \draw (opamp.corner 3) node[left, yshift=0.4cm] {$+$} node[left, yshift=-0.4cm] {$-$};
                % Entradas
                \draw (opamp.corner 3) ++(-0.5,-0.4) coordinate (inM);
                \draw (inM) -- ++(-0.5,0) coordinate (fork);
                \draw (fork) -- ++(0,0.6) to[R, l=$R_1$] ++(-1.5,0) node[left] {$V_1$};
                \draw (fork) -- ++(0,-0.6) to[R, l=$R_2$] ++(-1.5,0) node[left] {$V_2$};
                % Realimentacion
                \draw (inM) ++(0.2,0) -- ++(0,1.5) to[R, l=$R_f$] ++(1.8,0) -- ++(0,-1.5);
                \draw (opamp.corner 1) -- ++(0.5,0) node[right] {$V_{out}$};
                % No inversora a tierra
                \draw (opamp.corner 3) ++(0,0.4) -- ++(-0.5,0) -- ++(0,-0.3);
            \end{tikzpicture}
            \end{centering}
        \end{column}
    \end{columns}
\end{frame}

% Diapositiva 7: Amplificador Diferencial
\begin{frame}{5. Amplificador Diferencial (Restador)}
    \begin{block}{Ecuación General}
        Si $R_1 = R_3$ y $R_2 = R_f$:
        \[ V_{out} = \frac{R_2}{R_1} (V_2 - V_1) \]
    \end{block}
    \vspace{0.2cm}
    \textbf{Aplicaciones principales:}
    \begin{itemize}
        \item Amplifica únicamente la \textbf{diferencia} entre dos señales de entrada.
        \item Rechaza el ruido común presente en ambas líneas (Alto CMRR - Relación de Rechazo de Modo Común).
        \item Es la base de los amplificadores de instrumentación.
    \end{itemize}
\end{frame}

% Diapositiva 8: Integrador y Derivador
\begin{frame}{6. Configuraciones con Capacitores}
    \begin{columns}
        \begin{column}{0.5\textwidth}
            \begin{block}{Amplificador Integrador}
                Sustituye $R_f$ por un capacitor $C$.
                \[ V_{out} = -\frac{1}{R \cdot C} \int V_{in} \, dt \]
                Utilizado para generar rampas de voltaje y en filtros.
            \end{block}
        \end{column}
        \begin{column}{0.5\textwidth}
            \begin{block}{Amplificador Derivador}
                Sustituye $R_{in}$ por un capacitor $C$.
                \[ V_{out} = -R \cdot C \frac{dV_{in}}{dt} \]
                La salida es proporcional a la razón de cambio de la entrada.
            \end{block}
        \end{column}
    \end{columns}
\end{frame}

% Diapositiva 9: Conclusión
\begin{frame}{Resumen y Conclusiones}
    \begin{itemize}
        \item \textbf{Flexibilidad:} Cambiando los componentes de la red de realimentación externa ($R$ y $C$) se define por completo la función del circuito.
        \item \textbf{Realimentación Negativa:} Es la clave en todas estas configuraciones para estabilizar la ganancia y controlar el circuito de forma lineal.
        \item \textbf{Criterio de Selección:} La elección de la configuración depende de si necesitamos invertir la señal, sumar canales, aislar etapas o realizar operaciones matemáticas complejas.
    \end{itemize}
    \vspace{0.5cm}
    \centering
    \color{blue}\Large ¡Gracias por su atención!
\end{frame}

\end{document}